\documentclass[11pt]{article}
\usepackage[margin=2.6cm]{geometry}
\usepackage{amsmath,amssymb,amsthm}
\usepackage{xcolor}
\usepackage{enumitem}
\setlist{itemsep=1pt,topsep=3pt}

\newcommand{\Hom}{\mathrm{Hom}}
\newcommand{\I}{\mathbb{I}}
\newcommand{\C}{\mathcal{C}}
\newcommand{\V}{\mathsf{V}}
\newcommand{\U}{\mathsf{U}}
\newcommand{\El}{\mathsf{El}}
\newcommand{\Ext}{\mathrm{Ext}}
\newcommand{\ext}{\mathrm{ext}}
\newcommand{\glue}{\mathrm{glue}}
\newcommand{\Fib}{\mathrm{Fib}}
\newcommand{\HFib}{\mathrm{HFib}}
\newcommand{\full}{\mathrm{full}}
\newcommand{\Ty}{\mathrm{Ty}}
\newcommand{\STy}{\mathrm{STy}}
\newcommand{\Tm}{\mathrm{Tm}}
\newcommand{\RC}{\mathrm{RC}}
\newcommand{\Contr}{\mathrm{Contr}}
\newcommand{\Comp}{\mathrm{Comp}}
\newcommand{\Equiv}{\mathrm{Equiv}}
\newcommand{\Cl}{\mathrm{Cl}}
\newcommand{\GD}{\mathrm{GD}}
\newcommand{\SA}{\mathrm{SA}}
\newcommand{\Ua}{\mathrm{Ua}}
\newcommand{\ua}{\mathrm{ua}}
\newcommand{\isFib}{\mathrm{isFib}}
\newcommand{\Partial}{\mathrm{Partial}}
\newcommand{\Path}{\mathrm{Path}}
\newcommand{\Id}{\mathrm{Id}}
\newcommand{\id}{\mathrm{id}}
\newcommand{\up}{\mathrm{up}}
\newcommand{\code}[1]{\ulcorner #1\urcorner}
\newcommand{\Lam}{\Lambda}

\theoremstyle{plain}
\newtheorem{theorem}{Theorem}[section]
\newtheorem{lemma}[theorem]{Lemma}
\newtheorem{proposition}[theorem]{Proposition}
\newtheorem{corollary}[theorem]{Corollary}
\theoremstyle{definition}
\newtheorem{definition}[theorem]{Definition}
\newtheorem{construction}[theorem]{Construction}
\newtheorem{hypothesis}[theorem]{Hypothesis}
\theoremstyle{remark}
\newtheorem{remark}[theorem]{Remark}

\title{A model of univalence over a base category with an interval\\[3pt]
\large posets and simplicial sets as instances}
\author{}
\date{16 October 2026}

\begin{document}
\maketitle

\begin{abstract}
\noindent
We present a model of dependent type theory with a univalent universe of
fibrant types, over an abstract base category with an interval object and
a presheaf of cofibrations whose elements are decidable sieves. Fibrancy is
a structure: a uniform composition operation along lines. All the
structures of the model are instances of one notion, a relative
contractibility structure, and they are transported along maps of
contexts by one lemma. Following Voevodsky, every construction on a
family of data (the glued type of Theorem~\ref{thm:7}, its classifier, the
univalence operation) is made once on a universal family and obtained by
substitution along a classifying map; naturality is then the composition
law of substitution together with one equation of data. The two instances
are the category of finite non-empty posets, where the argument is
constructive, and the simplex category $\Delta$, where excluded middle is
used at exactly one point, the passage from horn filling to box filling.
Finally, we give identity types with a strict computation rule: Swan's,
for both instances, equivalent to the path type; and, on $\Delta$, the
path type itself, using excluded middle a second time.
\end{abstract}

\tableofcontents

%======================================================================
\section{Metatheory}\label{sec:meta}

We work in dependent type theory with $\Pi$, $\Sigma$, identity types,
natural numbers, Booleans, and a cumulative hierarchy of universes. We
assume function extensionality and univalence for the universes of the
metatheory; both may be used anywhere. (Function extensionality follows
from univalence; we use it directly, and the development is also obtained
from univalence alone by one instance, Appendix.) Logic is constructive, with one
exception: in Section~\ref{sec:delta}, and only there, we use excluded
middle for propositions,
\[
\mathrm{lem}:\ \Pi(P:\mathrm{Prop})\ P+\neg P .
\]
The axiom of choice is never used (Remark~\ref{rem:ac}).

\paragraph{Iterative sets.} For each universe level $a$ we use the type
$\V_a$ of iterative sets (Aczel's sets as well-founded trees, with
extensional equality; in univalent type theory, the type of
well-founded trees whose branching maps are injective). An element
$v:\V_a$ has a type of elements $\mathrm{Ix}(v)$, which is a set. Every
set in the model is \emph{coded}: the presheaves of the model are
families of iterative sets, and two types of the model are equal exactly
when their codes are equal in $\V$.

\begin{remark}[strictness]\label{rem:strict}
Equality in $\V_a$ is extensional, as in set theory, not univalent: two
isomorphic codes are in general different elements. The metatheory is
univalent, but this does not make isomorphic types of the model equal.
All uniformity conditions of the model (the uniformity of a composition,
the naturality of the univalence operation, the stability of $\Ext$ under
substitution) are \emph{equalities of codes}. Never an isomorphism, never
a relation.
\end{remark}

Each level $a$ gives a level $a^+$ with $\V_a$ coded in $\V_{a^+}$: a code
$v:\V_a$ has a code $\up(v):\V_{a^+}$ with $\mathrm{Ix}(\up(v))$ and
$\mathrm{Ix}(v)$ the same set. Types of the model live at level $a$; the
universe and the contexts built from it live at level $a^+$.

%======================================================================
\section{The base}\label{sec:base}

\subsection{The category and the interval}

The base is a small category $\C$ with an object $[1]$ and, for every
object $X$, two maps $e_0,e_1:X\to[1]$, natural in $X$ (the constant
maps at the end points). We write $\I=y[1]$ for the presheaf represented
by $[1]$, so $\I(W)=\Hom(W,[1])$. We do not assume that $\C$ has
products: posets and $\Delta$ are both instances, and $\Delta$ has no
product $X\times[1]$.

\paragraph{Connections.} There are two maps
$c_0,c_1:\I\times\I\to\I$ (on a stage $W$: for $i,j:W\to[1]$, a map
$c_k(i,j):W\to[1]$), natural in $W$, with
\[
c_k(e_k,j)=e_k,\qquad c_k(i,e_k)=e_k,\qquad
c_k(e_{1-k},j)=j,\qquad c_k(i,e_{1-k})=i .
\]
In both instances, for the order $e_0<e_1$, $c_0=\min$ and $c_1=\max$,
computed pointwise on maps $W\to[1]$.

\subsection{Presheaves, types, lines}

A \emph{context} is a presheaf $\Gamma$ on $\C$ with values in coded
sets: a family of codes $\Gamma(X)$ with restrictions $\rho\mapsto\rho f$.
A \emph{type} $A$ over $\Gamma$ is a family of codes $A(X,\rho)$ for
$\rho\in\Gamma(X)$ with restrictions $a\mapsto af:A(X,\rho)\to
A(W,\rho f)$, functorial. We write $E_A$ (or $\Gamma.A$) for the total
presheaf of the pairs $(\rho,a)$. A substitution $s:\Delta\to\Gamma$ is a
natural transformation, and $A[s](X,\rho)=A(X,s\rho)$.

\begin{definition}[lines]
A \emph{line} of $\Gamma$ at $X$ is a family $\gamma(W,f,j)\in\Gamma(W)$
for $f:W\to X$ and $j\in\I(W)$, natural:
$\gamma(W,f,j)\cdot g=\gamma(V,fg,jg)$. Equivalently a map
$yX\times\I\to\Gamma$. The lines form a presheaf $\Gamma^\I$, with
evaluations $\gamma\mapsto\gamma(e)=\gamma(X,\id,e)$ for $e=e_0,e_1$ and
constant lines $\rho\mapsto(W,f,j)\mapsto\rho f$. A line in $A$ over a
line $\gamma$ of $\Gamma$ is a line of $E_A$ lying over $\gamma$; these
form a type $A^\I$ over $\Gamma^\I$.
\end{definition}

This is the path object of a presheaf, defined without products in $\C$
(the exponential law holds by definition). On posets $yX\times\I=y(X\times[1])$,
so a line at $X$ is an element of $\Gamma(X\times[1])$.

\subsection{Cofibrations}

\begin{definition}[cofibrations]
For each object $X$ we are given a set $\Phi(X)$ with a map sending
$\varphi\in\Phi(X)$ to a \emph{decidable sieve} on $X$: a decidable
property ``$f$ lands in $\varphi$'' of maps $f:W\to X$, closed under
precomposition. Two elements with the same maps landing in them are
equal. $\Phi$ is a presheaf: $f^*\varphi$ is the sieve of the $g$ with
$fg$ in $\varphi$. We assume:
\begin{enumerate}
\item the empty sieve $\bot$ and the total sieve $\top$ are in $\Phi$;
  whether $\varphi=\top$ is decidable (it is whether $\id_X$ lands in
  $\varphi$);
\item $\Phi(X)$ is closed under finite unions;
\item the points: the sieves $(j=e_0)$ and $(j=e_1)$ on $W$, for
  $j:W\to[1]$, are in $\Phi(W)$;
\item $\forall$ along the interval: for a sieve $L$ of $yX\times\I$
  (a decidable family $L(W,f,j)$, closed under restriction) given by
  elements of $\Phi$ on the generalised points, there is
  $\forall L\in\Phi(X)$ with
  \[
  f\in\forall L\iff L(V,fg,j)\ \text{for all }g:V\to W,\ j\in\I(V).
  \]
\end{enumerate}
\end{definition}

Since an element of $\Phi$ \emph{is} a sieve, a cofibration given by
several faces is one element, their union; a list of faces is only a
notation for it.

\begin{remark}
On posets the subsets of $X$ are not closed under union as sieves: on
$[1]$, the union of the subsets $\{0\}$ and $\{1\}$ is the whole of
$[1]$ as a subset, but the identity of $[1]$ lands in neither. So the
sieves of the poset instance are the decidable sieves of maps, as on
$\Delta$.
\end{remark}

A \emph{cofibration} over $\Gamma$ is a natural family
$\psi_\rho\in\Phi(X)$ for $\rho\in\Gamma(X)$. The point $\rho$ is
\emph{full} if $\psi_\rho=\top$; this is decidable, and $f$ is full
($\rho f$ is full) exactly when $f$ lands in $\psi_\rho$. We write
$\Gamma|\psi$ for the context of full points.

%======================================================================
\section{Relative contractibility}\label{sec:rc}

\subsection{Partial elements}

\begin{definition}
Let $S$ be a type over $\Gamma$ and $\rho\in\Gamma(X)$. A \emph{partial
element} of $S$ at $\rho$ is a pair $(\varphi,u)$ with $\varphi\in\Phi(X)$
and $u$ a family $u_f\in S(W,\rho f)$ for $f:W\to X$ landing in
$\varphi$, natural: $u_f\cdot g=u_{fg}$. Its \emph{restriction} along
$h:Y\to X$ is $(h^*\varphi,\ g\mapsto u_{hg})$ at $\rho h$.
\end{definition}

Partial elements at $\rho$ form a presheaf $\Partial(S)$ over $\Gamma$
(cours6: $\tilde A$). Two partial elements are equal when their sieves
are equal and their values agree; the values do not depend on the proof
that $f$ lands.

\subsection{The notion}

\begin{definition}[relative contractibility structure]\label{def:rc}
Let $S$ and $T$ be types over $\Gamma$ and $m:S\to T$ a map over
$\Gamma$. A \emph{relative contractibility structure} $k\in\RC(\Gamma,S,T,m)$
assigns to each $\rho\in\Gamma(X)$ and each partial element $(\varphi,u)$
of $S$ at $\rho$ an element $k(\rho,\varphi,u)\in T(X,\rho)$ such that
\begin{enumerate}
\item (sides) $k(\rho,\varphi,u)\cdot f=m(u_f)$ for $f$ landing in
  $\varphi$;
\item (uniformity) $k(\rho,\varphi,u)\cdot h=k(\rho h,h^*(\varphi,u))$
  for $h:Y\to X$.
\end{enumerate}
\end{definition}

The two conditions are propositions, so:

\begin{lemma}[extensionality]\label{lem:rc-ext}
Two relative contractibility structures with the same values are equal.
\end{lemma}

\begin{remark}[property gives structure]
A structure $k$ is a map $\Partial(S)\to T$ over $\Gamma$ restricting to
$m$ on the sieve; uniformity is its naturality. It is one map out of the
generic partial element. No choice is involved.
\end{remark}

\subsection{Instances}

\begin{definition}
\leavevmode
\begin{enumerate}
\item A \emph{contractibility structure} on $S$ is
  $\Contr(\Gamma,S)=\RC(\Gamma,S,S,\id)$: every partial element extends
  to a total one, uniformly.
\item For $\alpha:B\to A$ over $\Gamma$, the \emph{strict fibre} of
  $\alpha$ at $(\rho,a)\in E_A$ is the set of $b\in B(\rho)$ with
  $\alpha b=a$; this is a type $\mathrm{SF}_\alpha$ over $E_A$. A
  \emph{trivial fibration structure} on $\alpha$ is
  $\Contr(E_A,\mathrm{SF}_\alpha)$.
\item \emph{Filling at $e$}: $\mathrm{Fill}_e(A)$ is a trivial fibration
  structure on the evaluation $\mathrm{ev}_e:A^\I\to A[\mathrm{ev}_e]$
  over $\Gamma^\I$: a partial line with a total end at $e$ extends to a
  total line with that end.
\item \emph{Composition from $e$ to $e'$}: let $L_e$ be the type, over the
  context $\Gamma^\I.A[\mathrm{ev}_e]$ of pairs (line $\gamma$, element
  $a$ over $\gamma(e)$), of the lines in $A$ over $\gamma$ ending at $a$
  at $e$, and $m=\mathrm{ev}_{e'}:L_e\to A[\mathrm{ev}_{e'}]$. Then
  \[
  \Comp_{e\to e'}(A)=\RC\bigl(\Gamma^\I.A[\mathrm{ev}_e],\ L_e,\
  A[\mathrm{ev}_{e'}],\ \mathrm{ev}_{e'}\bigr).
  \]
  A composition takes a partial line with a total start $a$ and gives an
  element at the other end that agrees with the ends of the sides.
\item For $\alpha:B\to A$ over $\Gamma$, the \emph{homotopy fibre}
  $\HFib_\alpha$ over $E_A$ has at $(\rho,a)$ the pairs $(b,\ell)$ with
  $b\in B(\rho)$ and $\ell$ a line in $A$ over the constant line at $\rho$
  with $\ell(e_0)=\alpha b$ and $\ell(e_1)=a$. An \emph{equivalence
  structure} on $\alpha$ is $\Equiv(\alpha)=\Contr(E_A,\HFib_\alpha)$.
\end{enumerate}
\end{definition}

\begin{remark}
Because a partial element is given on one sieve, a structure cannot
distinguish two presentations of the same sieve. No condition of
invariance under re-presentation (inserting an empty face, permuting
faces) is needed.
\end{remark}

\subsection{Moving a structure along a map}

\begin{construction}[move]\label{con:move}
Let $s:\Delta\to\Gamma$ be a map of contexts, possibly from a context of
level $a$ to one of level $a^+$. Let $m':S'\to T'$ over $\Delta$ and
$m:S\to T$ over $\Gamma$, and let $\sigma:S'\to S[s]$ and
$\tau:T'\to T[s]$ be maps over $\Delta$ with $\tau m'=m\sigma$, $\tau$
a bijection on each fibre. For $k\in\RC(\Gamma,S,T,m)$ define
\[
s^*k(\rho,\varphi,u)=\tau^{-1}\bigl(k(s\rho,\varphi,\sigma u)\bigr).
\]
\end{construction}

\begin{lemma}
$s^*k\in\RC(\Delta,S',T',m')$.
\end{lemma}

\begin{proof}
On a side $f$:
$s^*k(\ldots)\cdot f=\tau^{-1}(k(s\rho,\varphi,\sigma u)\cdot f)
=\tau^{-1}(m\sigma u_f)=\tau^{-1}\tau m'u_f=m'u_f$, using that $\tau$
and $\tau^{-1}$ commute with restriction. Uniformity follows in the same
way from the naturality of $s$, $\sigma$ and $\tau$.
\end{proof}

\begin{lemma}[laws]\label{lem:move-laws}
Moving along the identity with identity maps is the identity. Moving
along $s$ and then along $t$ is moving along $st$ with the composite
maps. Moving commutes with substitution, and is a congruence in its data.
\end{lemma}

\begin{proof}
In each case the two structures have the same values, by the definition
of $s^*k$; conclude by Lemma~\ref{lem:rc-ext}.
\end{proof}

\begin{lemma}[transport is a move]\label{lem:tr-move}
Let $P:A=A'$ be an equality of types over $\Delta$ and $c$ a fibration
structure on $A$. Then $(A,c)=(A',c')$, where $c'$ is $c$ moved along the
cartesian map whose fibre maps are transport along $P$. These fibre maps
may be taken to be transport along any equality of the fibres, since $V$
is a set.
\end{lemma}

\begin{proof}
By induction on $P$; at reflexivity the cartesian map is the identity,
and Lemma~\ref{lem:move-laws} gives $c'=c$.
\end{proof}

In practice the move is used at the level of types. A \emph{cartesian
map} $A'\to A$ over $s:\Delta\to\Gamma$ is a map $A'\to A[s]$ over
$\Delta$ which is a bijection on each fibre. It induces the maps $\sigma$
and $\tau$ for every structure of the model built from $A$ (for
compositions: on partial lines and on end points), so that a cartesian
map moves $\Comp$ and $\Fib$ in one step. All the transfers of the model
are instances of a cartesian map: substitution along $s$ (the identity
$A[s]\to A[s]$); a type of level $a$ read at level $a^+$ (the fibres are
the same sets); and a type isomorphic fibrewise to the substitution of
another.

%======================================================================
\section{Fibrations}\label{sec:fib}

\begin{definition}
A \emph{fibration structure} on $A$ over $\Gamma$ is a pair of
compositions
\[
\Fib(\Gamma,A)=\Comp_{e_0\to e_1}(A)\times\Comp_{e_1\to e_0}(A).
\]
A \emph{fibrant type} over $\Gamma$ is a pair $(A,c)$ with
$c\in\Fib(\Gamma,A)$.
\end{definition}

The composition is the stored structure; filling is derived.

\begin{proposition}[filling from composition]\label{prop:fill}
A composition $\Comp_{e\to e'}(A)$ gives $\mathrm{Fill}_e(A)$, uniformly.
The construction commutes with moves: the filling of a moved composition
is the moved filling (Construction~\ref{con:move}).
\end{proposition}

\begin{proof}
Given a partial line $(\varphi,u)$ in $A$ over $\gamma$ with total end
$a$ at $e$, consider at the stage $(W,f,i)$ the line
$j\mapsto\gamma(f,c(i,j))$ along the connection that is constant at $e$.
Its partial line has the sides of $u$ connected, and a side on $(i=e)$
given by $a$. Its composite at $e'$ is the value of the filler at $i$.
The connection laws give the end at $e$ and the sides; uniformity is that
of the composition. For moves: the filling problem at a point is formed
from the data at that point only, so the problem of the moved composition
is the moved problem; conclude by Lemma~\ref{lem:rc-ext}.
\end{proof}

\begin{lemma}[end of the filler]\label{lem:fill-end}
At $i=e'$ the filler of Proposition~\ref{prop:fill} is the composite of
the composition it is built from.
\end{lemma}

\begin{proof}
At $i=e'$ the extra side $(i=e)$ is empty, since $e$ and $e'$ are
disjoint, so the sieve of the problem is the original one; by the
connection laws the line $j\mapsto\gamma(f,c(e',j))$ is $\gamma f$
itself. The two problems are equal, hence so are their composites.
\end{proof}

\begin{proposition}[substitution]
For $s:\Delta\to\Gamma$, a fibration structure on $A$ moves to one on
$A[s]$ (Construction~\ref{con:move}, with lines moved along $s^\I$).
The laws $c[\id]=c$ and $c[st]=c[s][t]$ hold strictly
(Lemma~\ref{lem:move-laws}).
\end{proposition}

\begin{theorem}[type formers]\label{thm:formers}
Fibrant types are closed under $\Sigma$, $\Pi$ and path types, naturally
in $\Gamma$ (for $s:\Delta\to\Gamma$, the structure on $\Sigma(A[s],B[s])$
is the move of the one on $\Sigma(A,B)$ along the cartesian map
$\Sigma(A[s],B[s])\to\Sigma(A,B)$ over $s$, and likewise for $\Pi$ and
paths):
\begin{gather*}
\Fib(\Gamma,A)\to\Fib(\Gamma.A,B)\to\Fib(\Gamma,\Sigma_AB),\qquad
\Fib(\Gamma,A)\to\Fib(\Gamma.A,B)\to\Fib(\Gamma,\Pi_AB),\\
\Fib(\Gamma,A)\to\Fib(\Gamma,\Path_A(a,b)).
\end{gather*}
\end{theorem}

\begin{proof}
$\Sigma$: fill the first component in $A$, then the second in $B$ over
the line of $\Gamma.A$ so obtained (a line in $A$ over a line in $\Gamma$
is a line in $\Gamma.A$). $\Pi$: for a line of arguments ending at a
given argument, fill it backwards in $A$, apply the partial line of
functions to the filled line, and compose in $B$. $\Path$: a partial line
of paths is a partial line in $A$ over the lines $\gamma(f,j)$ with the
two end conditions as further sides on $(i=e_0)$ and $(i=e_1)$; these
three sides are joined into one sieve (their union), and one composes in
$A$. As for filling (Proposition~\ref{prop:fill}), the problem is formed
from the input restricted along the identity, so that one uniformity
argument gives both its naturality and its uniformity.

Naturality: each construction is formed pointwise from the structures
of $A$ and $B$ and from filling, which commutes with moves
(Proposition~\ref{prop:fill}); so the structure on the substituted type
and the moved one have the same values, and are equal by
Lemma~\ref{lem:rc-ext}. The form with a cartesian map avoids an equality
of types $\Sigma(A[s],B[s])=\Sigma(A,B)[s]$, which holds only up to an
identification of the actions; with it, Lemma~\ref{lem:tr-move} gives the
substitution form.
\end{proof}

\begin{proposition}[transport is an equivalence]\label{prop:transport}
Let $A$ be fibrant over $\Gamma$ and $\gamma$ a line of $\Gamma$ at $X$.
The end at $e'$ of the filler from $e$ is a map
$\mathrm{tr}_\gamma:A(\gamma(e))\to A(\gamma(e'))$ over $yX$, natural in
$\gamma$, with an equivalence structure, natural in $\gamma$.
\end{proposition}

\begin{proof}
The construction is made once over the universal line, over $\Gamma^\I$,
and substituted along $\gamma:yX\to\Gamma^\I$. The equivalence structure
extends a partial fibre element $(b,\ell)$ by filling, along the given
line, the paths of the fibre; this is a composition in a path type over
a line.
\end{proof}

%======================================================================
\section{Small types and the universe}\label{sec:univ}

\subsection{Small types over large contexts}

\begin{definition}
Let $G$ be a context of level $a^+$. A \emph{small type} over $G$ is a
type whose fibres are codes of level $a$: $A(X,x):\V_a$ for $x\in G(X)$.
We write $\STy(G)$; $\up A$ is $A$ read as a type of level $a^+$ (the
same fibres). For a context $D$ of level $a$ and a map $s:D\to G$, the
substitution $s^*A$ is the type of level $a$ with
$(s^*A)(X,\rho)=A(X,s\rho)$.
\end{definition}

\begin{lemma}[composition law]\label{lem:xsub-o}
For $t:T\to D$, $(s^*A)[t]=(st)^*A$, as an equality of types.
\end{lemma}

\begin{proof}
The fibres agree by definition; the restrictions agree because they
differ only by the proof of naturality used to transport along, and the
contexts are sets.
\end{proof}

This is Voevodsky's chosen pullback: substitution of small types is
strictly functorial, once for all small types. A structure on $\up A$
(a fibration structure, an equivalence structure) moves to $s^*A$ by
Construction~\ref{con:move}, with identity fibre maps, and its laws are
those of Lemma~\ref{lem:move-laws}.

\subsection{The universe}

\begin{definition}[classifier]
A \emph{fibrancy classifier} assigns to each context $\Gamma$ (of level
$a$) and type $A$ over $\Gamma$ a set $\Cl(\Gamma,A)$, with restrictions
$c\mapsto c[s]$ along $s:\Delta\to\Gamma$ satisfying
\[
(A,c)[\id]=(A,c),\qquad (A,c)[st]=(A,c)[s][t]
\]
as equalities in $\Sigma_A\Cl(\Gamma,A)$, and such that $\Cl(yX,A)$ has a
code in $\V_{a^+}$.
\end{definition}

\begin{definition}[the universe]
$\U(X)=\Sigma(A\in\Ty_a(yX))\,\Cl(yX,A)$, with restriction along
$f:Y\to X$ given by substitution along $yf$. This is a presheaf of codes
of level $a^+$. $\El\in\STy(\U)$ is the small type with
$\El(X,(A,c))=A(X,\id_X)$.
\end{definition}

A type $A$ over $\Gamma$ with $c\in\Cl(\Gamma,A)$ has a \emph{code map}
$\code{A}:\Gamma\to\U$, $\rho\mapsto(A[\hat\rho],c[\hat\rho])$, where
$\hat\rho:yX\to\Gamma$ is the Yoneda map of $\rho$; and
$\code{A}^*\El=A$ as types. A small type over a context of codes is read
by $\El$: for any $c:\Delta\to\U$, $c^*\El$ is ``El along $c$''.

\begin{hypothesis}[El is fibrant]\label{hyp:el}
$\up\El$ has a fibration structure over $\U$.
\end{hypothesis}

Then every type of the form $c^*\El$ is fibrant, by moving $\El$'s
structure along $c$; the structures so obtained are natural in $c$ by
Lemma~\ref{lem:move-laws}. Every fibration structure used by the
constructions below is of this form.

\begin{remark}
The classifier is a parameter because the two instances differ:
on posets $\Cl=\Fib$ and Hypothesis~\ref{hyp:el} holds because $\I$ is
tiny (Section~\ref{sec:posets}); on $\Delta$, $\I$ is not tiny, $\Cl$ is
horn filling, and Hypothesis~\ref{hyp:el} uses excluded middle
(Section~\ref{sec:delta}).
\end{remark}

%======================================================================
\section{Gluing: Theorem 7}\label{sec:glue}

\subsection{$\Ext(\alpha)$}

\begin{definition}
Let $A$ be a type over $\Gamma$, $\psi$ a cofibration over $\Gamma$, $B$ a
type over $\Gamma|\psi$ and $\alpha:B\to A|\psi$ a map. For
$\rho\in\Gamma(X)$:
\begin{itemize}
\item if $\rho$ is full, $\Ext(\alpha)(X,\rho)=B(X,\rho)$;
\item otherwise $\Ext(\alpha)(X,\rho)$ is the set of pairs $(a,b)$ with
  $a\in A(X,\rho)$ and $b$ a natural family $b_f\in B(W,\rho f)$ for the
  full $f:W\to X$, with $\alpha(b_f)=af$.
\end{itemize}
Restriction is by cases on the source and the target (a full point
restricts to a full point). The map $\ext:\Ext(\alpha)\to A$ is $\alpha$
on full points and the first projection otherwise. The map $\glue$ sends
a pair $(a,b)$ to $b_\id$ at a full point and to $(a,b)$ otherwise; so
$\ext(\glue(a,b))=a$ in both cases.
\end{definition}

On $\psi$, $\Ext(\alpha)$ \emph{is} $B$ and $\ext$ \emph{is} $\alpha$: as
codes, not up to isomorphism.

\begin{theorem}[Theorem 7]\label{thm:7}
Let $A$ and $B$ be fibrant and $e_\alpha\in\Equiv(\alpha)$. Then
$\Ext(\alpha)$ has a fibration structure and $\ext$ an equivalence
structure, naturally in $\Gamma$. On full points the composition of
$\Ext(\alpha)$ is that of $B$ and the equivalence structure of $\ext$ is
$e_\alpha$.
\end{theorem}

\begin{proof}[Proof: $\ext$ is an equivalence]
Let $w=(\rho,a)\in E_A(X)$ and $(\varphi,p)$ a partial element of
$\HFib_\ext$ at $w$.
\begin{enumerate}
\item For full $f:W\to X$, $p$ restricted along $f$ is a partial element
  of $\HFib_\alpha$ at $wf$, since $\Ext(\alpha)$ is $B$ there; let
  $\theta_f\in\HFib_\alpha(wf)$ be its extension by $e_\alpha$. Then
  $\theta$ is natural (uniformity of $e_\alpha$), equals $p_f$ on
  $\varphi$ (sides of $e_\alpha$, with the identity in the restricted
  sieve), and does not depend on the proof that $f$ is full.
\item Consider the partial line in $A$ over the constant line at $\rho$,
  with total end $w$ at $e_1$, with sides: the lines of $p$ on $\varphi$,
  and the images in $A$ of the lines of $\theta$ on $\psi_\rho$. They
  agree where both are defined, by (1). Fill it by
  $\mathrm{Fill}_{e_1}(A)$; call the filler $F$.
\item Let $a_0=F(e_0)$ and $b_f$ the $B$ component of $\theta_f$. Then
  $\alpha(b_f)=$ the line of $\theta_f$ at $e_0$ $=F(f,e_0)=a_0f$, so
  $(a_0,b)$ is a pair, and $v_0=(\glue(a_0,b),F)\in\HFib_\ext(w)$.
\end{enumerate}
$v_0$ restricts to $p$ on $\varphi$ by (1) and the sides of the filler,
and is uniform because every step is. At a full $\rho$, $\psi_\rho=\top$,
so $F$ is the line of $\theta_\id$, which is the extension by $e_\alpha$.
No $\forall$ is needed: the line of $\Gamma$ is constant.
\end{proof}

\begin{proof}[Proof: composition]
Let $\gamma$ be a line of $\Gamma$ at $X$, $x$ an element of
$\Ext(\alpha)$ over $\gamma(e)$, and $(\varphi,u)$ a partial line in
$\Ext(\alpha)$ over $\gamma$ ending at $x$. Let $\delta=\forall_\gamma\psi$:
the maps $f$ such that $\gamma(fg,j)$ is full for all $(g,j)$.
\begin{enumerate}
\item On $\delta$ the problem is a problem of $B$ over a line of
  $\Gamma|\psi$; fill it in $B$: a partial line $\beta$ in $B$ on
  $\delta$.
\item In $A$: the partial line with sides $\ext\circ u$ on $\varphi$ and
  $\alpha\circ\beta$ on $\delta$, total end $\ext(x)$ at $e$, is filled;
  let $a_1$ be its end at $e'$.
\item At $\gamma(e')$, on $\psi_{\gamma(e')}$, the ends of $u$ (on
  $\varphi$) and of $\beta$ (on $\delta$) give a partial element of the
  homotopy fibre of $\alpha$ at $a_1$ (with constant paths); extend it by
  $e_\alpha$, on the sieve $\psi_{\gamma(e')}$, to $(b_1,\omega)$.
\item Fill in $A$ the line $\omega$ backwards from $a_1$, with sides
  $\ext$ of the ends of $u$ on $\varphi$, of $\beta$ on $\delta$ (constant
  lines), and $\omega$ on $\psi_{\gamma(e')}$; its end $a_2$ satisfies
  $\alpha(b_{1})=a_2$ on $\psi_{\gamma(e')}$.
\item The composite is $\glue(a_2,b_1)$.
\end{enumerate}
It agrees with the sides of $u$ by construction, and is uniform because
$\forall$, the fillers and $e_\alpha$ are. If $\gamma$ is full
everywhere, $\delta=\top$ and the composite is $B$'s.
\end{proof}

\subsection{The universal glue data}

\begin{definition}[glue data]
The presheaf $\GD$ of level $a^+$ has as points at stage $I$ the tuples
$d=(A,\varphi,B,\alpha,e)$ where
\begin{itemize}
\item $A\in\U(I)$ and $\varphi\in\Phi(I)$;
\item $B$ assigns to each $f:W\to I$ landing in $\varphi$ a code
  $B_f\in\U(W)$, natural;
\item $\alpha$ is a map $\El(B)\to\El(A)$ over the part of $yI$ landing in
  $\varphi$, natural;
\item $e$ is an equivalence structure on $\alpha$ (for $\El$'s
  fibrations).
\end{itemize}
Restriction along $h:J\to I$ is precomposition; its laws are equalities
of codes.
\end{definition}

The restriction of the component $e$ along $h$ is its move
(Construction~\ref{con:move}) along the cartesian map of triples
$\chi_h$ from $(A,B,\alpha)\cdot h$ to $(A,B,\alpha)$ over $yh$, whose
fibre maps are identities; the laws of restriction are those of
Lemma~\ref{lem:move-laws}.

A structure over a presheaf can be given locally, point by point.

\begin{lemma}[gluing structures along points]\label{lem:glue-rc}
Let $m:S\to T$ over a presheaf $D$. Suppose that for each $x\in D(Z)$
we are given a context $C_x$ with a map $\iota_x:C_x\to D$, a point
$o_x\in C_x(Z)$ with $\iota_x(o_x)=x$, and a structure
$k_x\in\RC(C_x,S[\iota_x],T[\iota_x],m[\iota_x])$. Suppose they are
coherent: for every $c\in C_x(Y)$ and every partial element $u$ at
$\iota_x(c)$, the extension given by $k_x$ at $c$ and the one given by
$k_{\iota_x(c)}$ at $o_{\iota_x(c)}$ have the same image in $T$ (in
particular for $c=o_x\cdot g$). Then
\[
k(x,\varphi,u)=k_x(o_x,\varphi,u)
\]
is a structure $k\in\RC(D,S,T,m)$, the unique one agreeing with every
$k_x$ at $o_x$; moreover $k$ moved along $\iota_x$ is $k_x$.
\end{lemma}

\begin{proof}
The sides of $k$ are those of the $k_x$. Uniformity along $g$ is the
uniformity of $k_x$ from $o_x$ to $o_x\cdot g$ followed by the coherence.
Uniqueness holds by Lemma~\ref{lem:rc-ext}.
\end{proof}

The case $C_x=yZ$, $o_x=\mathrm{id}_Z$ says that a structure is
determined by its values at the generic points. Stating the coherence in
$T$ means that no transport between the types $S[\iota_x][g]$ and
$S[\iota_{x\cdot g}]$ is needed.

Over $\GD$ there is the \emph{universal glue datum}: $A_\U$, $\psi_\U$,
$B_\U$, $\alpha_\U$, $e_\U$, where $A_\U$ and $B_\U$ are $\El$ along the
components $A$ and $B$, and $\psi_\U(d)=\varphi_d$. The structure $e_\U$
is glued by Lemma~\ref{lem:glue-rc} from the points' own structures:
for $d\in\GD(Z)$, $C_d$ is the context of the equivalence problems of
$\alpha_d$, $\iota_d$ is the cartesian map $(f,q)\mapsto(d\cdot f,q)$ to
the universal one, and $k_d=e_d$; the coherence is the definition of
restriction in $\GD$. Theorem~\ref{thm:7}
at level $a^+$ gives $\Ext_\U=\Ext(\alpha_\U)$ with its fibration and the
equivalence structure of $\ext_\U$, \emph{once}.

\begin{construction}[the small copy]
$\SA\in\STy(\GD)$ has at $d\in\GD(I)$ the fibre at the identity of
$\Ext(\alpha_d)$ over $yI$, where $\alpha_d$ is the datum $d$ read over
$yI$. It is fibrewise isomorphic to $\Ext_\U$ (the extents agree: $f$ is
full for $d$ iff it lands in $\varphi_d$), so it carries $\Ext_\U$'s
fibration structure, map to $A_\U$ and equivalence structure, moved by
Construction~\ref{con:move}.
\end{construction}

As in Voevodsky's universe categories, glue data over a context are
maps into the presheaf of glue data.

\begin{definition}[glue data over a context]
Glue data over $\Gamma$ are a map $\ell:\Gamma\to\GD$; substitution along
$s:\Delta\to\Gamma$ is composition, $\ell[s]=\ell\circ s$.
\end{definition}

\begin{remark}
By Yoneda, such an $\ell$ is the same as componentwise data over $\Gamma$:
a code map $A:\Gamma\to\U$, a cofibration $\psi$, a code map $B$ over the
extent, a map $\alpha$ of their $\El$'s and an equivalence structure $e$
(at $\rho\in\Gamma(X)$, $\ell\rho$ is the datum of these along $\rho$,
read over $yX$). This correspondence is not needed below.
\end{remark}

\begin{definition}\label{def:ext-ell}
$\Ext(\ell)=\ell^*\SA$, with the fibration, the map $\ext$ to $A$ and the
equivalence structure of $\SA$ moved along $\ell$
(Construction~\ref{con:move}). Theorem~\ref{thm:7} is thereby proved once,
at the universal data, and holds for all glue data by moving.
\end{definition}

\begin{lemma}[comparison]\label{lem:ext-cmp}
The fibre of $\ell^*\SA$ at $\rho$ is the fibre of the direct definition
of Section~\ref{sec:glue} for the datum $\ell\rho$: $B_{\ell\rho}(\mathrm{id})$
if the identity is in $\varphi_{\ell\rho}$, and the glued fibre otherwise.
\end{lemma}

\begin{proof}
The fibre is that of $\SA$ at $\ell\rho$, the fibre at the identity of
$\Ext$ of the datum $\ell\rho$ over $yX$. Conclude by the cases on
fullness.
\end{proof}

\begin{theorem}[naturality of Theorem~\ref{thm:7}]\label{thm:7-nat}
For $s:\Delta\to\Gamma$, there are cartesian maps
$\Ext(\ell\circ s)\to\Ext(\ell)$ and $A(\ell\circ s)\to A(\ell)$ over $s$,
with identity fibre maps, commuting with $\ext$; along them the
fibration and the equivalence structure of $\Ext(\ell\circ s)$ are the
moves of those of $\Ext(\ell)$. (With Lemma~\ref{lem:tr-move} this gives
the substitution form $\Ext(\ell)[s]=\Ext(\ell\circ s)$, as for the type
formers.)
\end{theorem}

\begin{proof}
By Definition~\ref{def:ext-ell} the structures of $\Ext(\ell)$ and
$\Ext(\ell s)$ are moves of those of $\SA$ along $\ell$ and $\ell s$; the
cartesian maps are those of Lemma~\ref{lem:xsub-o}, and the composite and
congruence laws of Lemma~\ref{lem:move-laws} give the equalities.
\end{proof}

%======================================================================
\section{Univalence and the fibrancy of the universe}\label{sec:ua}

\subsection{The univalence operation}

\begin{definition}
Over $\U$, $\Ua$ is the type with, at $F\in\U(X)$, the triples
$(B,\alpha,e)$: $B\in\U(X)$, $\alpha:\El(B)\to\El(F)$ over $yX$, and
$e\in\Equiv(\alpha)$. Restriction along $f$ is substitution along $yf$
on $B$ and $\alpha$; on $e$ it is the move (Construction~\ref{con:move})
along the cartesian map of triples over $yf$, as for $\GD$. Its laws are
equalities of codes for $B$ and $\alpha$, and Lemma~\ref{lem:move-laws}
for $e$. $\Ua(F)$ is $\Sigma(Y:\U)\,\Equiv(Y,F)$.

Equivalently, by Yoneda, $\Ua(F)$ at $X$ is the set of glue data
$d\in\GD(X)$ with $A_d=F$ and $\varphi_d=\top$, and restriction is that
of $\GD$; its laws are then those of $\GD$.
\end{definition}

\begin{theorem}[univalence]\label{thm:ua}
There is $\ua\in\Contr(\U,\Ua)$, uniform as an equality of codes.
\end{theorem}

The univalence operation is obtained from one universal datum.

\begin{construction}
The \emph{universal univalence datum} $\mathcal U$ over $\GD$ is
$(\SA,\ \ext,\ \mathrm{extEquiv})$: the small copy of $\Ext_\U$ with
$\ext_\U$ and its equivalence structure moved onto it, with codomain
$A_\U$. The code of $\SA$ along a map $s:\Delta\to\GD$ is given by a
classifier element $\mathrm{smallCl}(s)\in\Cl(\Delta,s^*\SA)$, the move
along $s$ of one universal element of $\Cl(\GD,\SA)$, with
\begin{enumerate}
\item[(C1)] for $t:\Delta'\to\Delta$, the equality
  $\bigl((s^*\SA)[t],\ \mathrm{smallCl}(s)[t]\bigr)
  =\bigl((st)^*\SA,\ \mathrm{smallCl}(st)\bigr)$
  in $\Sigma_{A}\Cl(\Delta',A)$, whose first component is
  Lemma~\ref{lem:xsub-o}; equivalently, $\mathrm{smallCl}(s)[t]$ is the
  transport of $\mathrm{smallCl}(st)$ along Lemma~\ref{lem:xsub-o}. It is
  the composite law of Lemma~\ref{lem:move-laws};
\item[(C2)] for $y\in\GD(W)$ full at the identity,
  $\bigl(\hat y^*\SA,\ \mathrm{smallCl}(\hat y)\bigr)=B_y(\mathrm{id})$
  as elements of $\U(W)$.
\end{enumerate}

When the classifier is $\Fib$ (as for posets), (C2) is proved once from two inputs
of the instance: the fibrancy of $\El$ and Lemma~\ref{lem:el-yon}. The
type equality $\hat y^*\SA=\El(B_y)$ is not assumed but derived. On full
points there is a cartesian map from $\El(B_y)$ into $\Ext_\U$ (the
Yoneda map followed by the inclusion of $B$ on full points), and the
fibres of $\hat y^*\SA$ are related to those of $\El(B_y)$ by a path
whose transport, read in $\Ext_\U$, is that cartesian map. Since the
cartesian map into $\Ext_\U$ is injective, this transport is natural,
which gives the equality of types. The fibrations then agree by
Lemma~\ref{lem:tr-move}, Lemma~\ref{lem:full-move}(iii),
Lemma~\ref{lem:move-laws} and Lemma~\ref{lem:el-yon}.
\end{construction}

Both steps of the proof below produce a full glue datum whose equivalence
structure is a given structure moved along cartesian maps. Such a datum
is determined by those maps.

\begin{lemma}[uniqueness]\label{lem:glue-unique}
Let $x\in\GD(X)$, and suppose its equivalence structure is the move of
a structure $\varepsilon$ along cartesian maps $t_A$, $t_B$ that agree
pointwise with the ones used to build a glue datum $y$ from
$\varepsilon$. If $x$ and $y$ have the same $A$, the same $\varphi$ and
the same codes $B$, then $x=y$.
\end{lemma}

\begin{proof}
The maps $\alpha$ need no hypothesis: on each side both are read off
the cartesian maps, by $t_A\circ\alpha=\ext\circ t_B$, so they agree
once $t_A$ and $t_B$ agree pointwise. The two equivalence structures are
moves of $\varepsilon$ along pointwise equal maps, hence equal by
Lemma~\ref{lem:move-laws}.
\end{proof}

It is used with three facts about full points.

\begin{lemma}[full points]\label{lem:full-move}
\begin{enumerate}
\item[(i)] For every $d\in\GD(Z)$, $e_d$ is $e_\U$ moved along the
  cartesian map of $d$'s point into the universal datum.
\item[(ii)] At a full point, the equivalence structure of $\ext_\U$ moved
  along a cartesian map equals $e_\U$ moved along the composite of that
  map with the identification $\Ext=B$.
\item[(iii)] On the full points, $B$ maps into $\Ext(\alpha)$ by a
  cartesian map, and $\Ext$'s fibration moved along it is $B$'s.
\end{enumerate}
\end{lemma}

\begin{proof}
(i) is the last clause of Lemma~\ref{lem:glue-rc}, which uses the
coherence at every local point, not only at $o_x\cdot g$. (ii) and (iii) are the full-point clauses of
Theorem~\ref{thm:7}, carried by Construction~\ref{con:move}: a moved
structure at $\rho$ is $\tau^{-1}(k(s\rho,\varphi,\sigma u))$, and $k$
at the full point $s\rho$ is given by those clauses; the remaining
identification is uniqueness of moves (Lemma~\ref{lem:move-laws}). For
(iii): on a full line the composite of $\Ext$ is the end of $B$'s filler
(Theorem~\ref{thm:7}), which is $B$'s composite by
Lemma~\ref{lem:fill-end}.
\end{proof}

\begin{proof}[Proof of Theorem~\ref{thm:ua}]
Let $F\in\U(X)$ and $(\varphi,u)$ a partial element of $\Ua$ at $F$,
$u_f=(B_f,\alpha_f,e_f)$. The proof has two steps.

\emph{(a) The datum of a partial element.} There is a unique
$d=(F,\varphi,B,\alpha,e)\in\GD(X)$ with $d\cdot f=u_f$ for $f$ in
$\varphi$. Its components $B$ and $\alpha$ are read off the $u_f$ at the
identity. Its equivalence structure is $e_\U$ along
$\ell_u:\varphi\to\GD$, $f\mapsto u_f$ (a map into the full points, the
sieve $\varphi$ read as a subpresheaf of $yX$), moved along the
cartesian map of triples. So no new gluing is needed: the $e_f$ were
glued once, in $e_\U$ (Lemma~\ref{lem:glue-rc}), and the uniformity of
$u$ is what makes $\ell_u$ a map of presheaves.

\emph{(b) The extension.} Let $\hat d:yX\to\GD$ and let
$\ua(F,\varphi,u)$ be the full glue datum given by $\mathcal U$ along
$\hat d$:
\[
\ua(F,\varphi,u)=\hat d^*\mathcal U\ \text{at the identity},
\]
with code $\mathrm{smallCl}(\hat d)$ (at $g$, the code of
$(d\cdot g)^*\SA$).

Steps (a) and (b) are two instances of one construction: a glue datum
read through cartesian maps over a map of points into $\GD$. Given $B$,
a target triple with its equivalence structure, and cartesian maps on
the $A$ and $B$ sides, it yields $\alpha$, the naturality of $\alpha$,
and $e$. Its map is $\ext$ followed by the
Yoneda isomorphism $\hat d^*A_\U\cong\El(F)$: the fibre of $\hat d^*A_\U$
at $f$ is $\El(F\cdot(f\circ\mathrm{id}))$, equal to $\El(Ff)$ only
up to $f\circ\mathrm{id}=f$, so this is a cartesian map whose fibre maps
are transports; the equivalence structure is moved along it.
\emph{Uniformity}: $h^*(\varphi,u)$ gives $dh$, and
$\widehat{dh}=\hat d\circ yh$; conclude by Lemma~\ref{lem:xsub-o},
Lemma~\ref{lem:move-laws} and (C1).
\emph{Sides}: for step (a), $d\cdot f=u_f$ on $\varphi$ by
Lemma~\ref{lem:glue-unique} and Lemma~\ref{lem:full-move}(i). For step
(b), by uniformity it suffices to treat a side containing the identity;
then $\varphi=\top$ and $d=u_\id$. The value of (b) has code $B_\id$ by
(C2), map $\alpha_\id$ (the Yoneda isomorphism is the identity at the
identity), and its equivalence structure is $e_\id$ by
Lemma~\ref{lem:glue-unique} with Lemma~\ref{lem:full-move}(ii) and (i).
Here the cartesian map of $u_\id$ into $\SA$ is the point map, then the
identification $\Ext=B$ at full points, then the inverse of the
comparison of $\SA$ with $\Ext_\U$; the inverse followed by the
comparison is the identity by Lemma~\ref{lem:move-laws}. The pointwise
check on the $B$ side uses only that the comparison is injective and
that $\varphi$ holds at full points. So the value is $u_\id$.
\end{proof}

\subsection{The universe is fibrant}

\begin{theorem}\label{thm:fibU}
For every context $\Gamma$, the constant type $\U$ over $\Gamma$ has a
fibration structure.
\end{theorem}

\begin{proof}
Composition from $e$ to $e'$: let $\gamma$ be a line of $\Gamma$, $F$ a
code at the start, and $(\varphi,\lambda)$ a partial line of codes, with
$\lambda_f(e)=Ff$. For each side, transport along $\lambda_f$ from $e'$
to $e$ (Proposition~\ref{prop:transport}, for $\El$ along $\lambda_f$)
is an equivalence $\El(\lambda_f(e'))\to\El(Ff)$, natural in $f$. This
is a partial element of $\Ua$ at $F$: concretely, the full glue datum
read through cartesian maps over the map $f\mapsto\lambda_f$ into the
lines of codes (the construction of step (a) of the proof of
Theorem~\ref{thm:ua}), so that its naturality is that of the
construction. Its extension by $\ua$ has a code
$Y\in\U(X)$, which is the composite. On a side, $Yf=\lambda_f(e')$ by
the sides of $\ua$; uniformity is that of $\ua$ and of transport. The
equivalence structures must be natural in $f$ strictly; this is where
Proposition~\ref{prop:transport}, and through it the fibrancy of path
types over a line (Theorem~\ref{thm:formers}), is used.
\end{proof}

\begin{remark}
This is cours6's route: the equivalence extension property gives
univalence ($\Sigma(Y:\U)\,\Equiv(Y,F)$ is contractible), and univalence
gives $\Comp(1,\U)$. No invariance under re-presentation of the
cofibration is needed, because a partial element is on one sieve.
\end{remark}

%======================================================================
\section{Instances}\label{sec:inst}

\subsection{Posets}\label{sec:posets}

$\C$ is the category of finite non-empty posets and monotone maps (or the
full subcategory on finite products of $[1]$), with $[1]=\{0<1\}$,
connections $\max$ and $\min$, and $\Phi(X)$ the decidable sieves on $X$.
Finite unions are pointwise; $\forall$ along the interval is decided at
the generic point $(\pi_1,\pi_2):W\times[1]\to W\times[1]$, since
$W\times[1]$ is an object.

The classifier is $\Cl(\Gamma,A)=\Fib(\Gamma,A)$, with the substitution of
Section~\ref{sec:fib}; it is coded because a composition is determined by
its composites. Hypothesis~\ref{hyp:el} holds because $\I$ is tiny:
$yX\times\I=y(X\times[1])$, so a line of $\U$ at $X$ is a code over
$y(X\times[1])$, with its own composition. The composition of $\El$ is
glued by Lemma~\ref{lem:glue-rc} from these: at a point $x$, the line's
code at the generic point of $X\times[1]$, read in $\El$ through the
Yoneda cartesian map; the coherence is that a code restricted along
$g\times[1]$ is the code at $xg$, together with Lemma~\ref{lem:el-yon}.

\begin{lemma}[$\El$ along a Yoneda code]\label{lem:el-yon}
Let $F\in\U(X)$, with its fibration $c_F$ (the classifier part of the
code). Read $\El$ along the Yoneda map $\hat F:yX\to\U$; its fibre at
$f$ is $\El(F\cdot(f\circ\mathrm{id}))$, which is related to $F$ by the
Yoneda cartesian map (fibre maps: transport along $f\circ\mathrm{id}=f$).
Then the fibration of $\El$ moved along this cartesian map is $c_F$.
\end{lemma}

\begin{proof}
By Lemma~\ref{lem:rc-ext} it suffices to compare composites. At a point
$f$ and a line over it, the glued composition of $\El$ is, by
definition, the composition of the line's code at the generic point of
$X\times[1]$; for a line of $\hat F$ this code is $F$ restricted along
$f\times[1]$, whose composition is $c_F$ restricted, because the
classifier is moved by Construction~\ref{con:move}. The two composites
then agree through the cartesian map, by Lemma~\ref{lem:move-laws}.
\end{proof} The classifier is moved by
Construction~\ref{con:move}, so (C1) is Lemma~\ref{lem:move-laws}. Condition (C2)
holds by Theorem~\ref{thm:7}: at a full point every line over $yX$ is
full, and there $\Ext$'s composition is $B$'s; this full-point
composition clause is carried to $\SA$ by Lemma~\ref{lem:full-move}(iii).
Everything is constructive.

\subsection{Simplicial sets}\label{sec:delta}

$\C=\Delta$, $[1]$ the 1-simplex, connections $\max$ and $\min$, $\Phi([n])$
the decidable sieves on $[n]$. Here $yX\times\I$ is not representable.
$\forall$ along the interval is decided on the shuffles: the prism
$[m]\times\Delta[1]$ is the union of the $m+1$ simplices
$P_k:[m+1]\to[m]\times[1]$, and every simplex of the prism factors
through one of them.

\begin{definition}[horn filling]
For $0\le k\le m+1$ let $\Lam^k[m+1]\in\Phi([m+1])$ be the horn. A
\emph{horn filling operation} for $A$ over $\Gamma$, $\isFib(\Gamma,A)$,
assigns to each $\rho\in\Gamma([m+1])$ and each partial element of $A$ at
$\rho$ on $\Lam^k[m+1]$ an extension. No uniformity is required.
\end{definition}

$\isFib$ is a structure, untruncated: an element is an operation.
$\Cl(\Gamma,A)=\isFib(\Gamma,A)$ is a classifier: substitution moves a
horn of $A[s]$ to a horn of $A$, fills it there and reads the filler back
(Construction~\ref{con:move} for operations without uniformity).

\begin{theorem}\label{thm:hornbox}
\leavevmode
\begin{enumerate}
\item (constructive) A fibration structure gives a horn filling
  operation.
\item (with excluded middle) A horn filling operation gives a fibration
  structure.
\end{enumerate}
\end{theorem}

\begin{proof}
(1) A horn is a retract of the prism along a retraction line
(Gabriel--Zisman): a family $r(f,i)$ with $r(f,e')=f$, $r(f,e)$ landing
in one face of the horn, and $r$ preserving the other faces. The
composition along $r$ extends the horn.

(2) First, \emph{face-uniform} fillers, constructively: a filling problem
at $[N]$ is solved by filling the shuffles $P_N,\ldots,P_0$ (or
$P_0,\ldots,P_N$) one after the other, each by one horn filling; the
result agrees with the sides and is uniform along face maps. Second,
\emph{normalised} fillers: a problem $b$ is \emph{degenerate} along
$s_i$ when $b=s_i^*d_i^*b$. This is a proposition, decided by excluded
middle; the normalised filler of a degenerate problem is
$s_i^*$ of the filler of $d_i^*b$ for the first such $i$, and the
face-uniform filler otherwise. By the Eilenberg--Zilber factorisation,
by induction, normalised fillers are uniform along every map. They form
a filling structure at each end, whose ends give the compositions.
\end{proof}

Excluded middle is used only in deciding degeneracy. The passage (2) is
not natural in $\Gamma$ (the first degeneracy is not preserved by
substitutions identifying points), so the constructions never compose
(2) with substitution: Hypothesis~\ref{hyp:el} is (2) applied once, to
$\El$, which has a horn filling operation by the classifier itself (a
horn of $\El$ at a code $(A,c)$ is a horn of $A$, filled by $c$).

\paragraph{The classifier of $\Ext$.} The code of $\SA$ is given by one
horn filling operation on $\Ext_\U$, at level $a^+$, \emph{by cases on
fullness}: at a full point, the horn operation of $B$ (read in $B$, since
$\Ext$ is $B$ there); otherwise, the horn operation of $\Ext_\U$'s
composition given by Theorem~\ref{thm:hornbox}(1). The case split is
pointwise, and this is legitimate because a horn filling operation
carries no uniformity: the universal element needs no naturality. Then
$\mathrm{smallCl}(s)$ is this operation pulled back along $s$
(Construction~\ref{con:move} for operations); (C1) is the composition
law of moving operations. (C2) is proved by the same chain as for
posets, with operations in place of fibrations: along the derived
equality $\hat y^*\SA=\El(B_y)$ (whose derivation does not use the
classifier), Lemma~\ref{lem:tr-move} and the move laws for operations;
at full points the universal operation is $B$'s by the case split (the
analogue of Lemma~\ref{lem:full-move}(iii)); and the analogue of
Lemma~\ref{lem:el-yon}: the operation of $\El$ moved along the Yoneda
cartesian map of a code $F$ is $F$'s own. Its proof is definitional: at
$\rho$, the operation of $\El$ at the code $F\cdot\rho$ is $F$'s at
$\rho\circ\mathrm{id}$, since the action on codes is substitution and
the operation of a substitution is the pulled back one. The passage
Theorem~\ref{thm:hornbox}(2)$\circ$(1) is never used, and excluded
middle reaches the final statements only through
Hypothesis~\ref{hyp:el}, i.e.\ through (2) applied once, to $\El$. The univalence operation and
Theorem~\ref{thm:fibU} follow as in Section~\ref{sec:ua}.

\begin{remark}[choice]\label{rem:ac}
The usual Kan condition is a \emph{property}: every horn merely has a
filler. To obtain from it a horn filling \emph{operation} requires the
axiom of choice, and in univalent foundations choice yields only the
mere existence of such an operation. The model uses the operation
$\isFib$ as its definition and never passes through the property. How
faithfully this represents the usual mathematics of Kan complexes, where
filling is a property, is left for further analysis.
\end{remark}

%======================================================================
\section{Identity types}\label{sec:id}

The path type of Theorem~\ref{thm:formers} has an eliminator $J$ defined
by composition, but its computation rule holds only up to a path. We give
two identity types with a strict computation rule: Swan's, for any base
with one more property of $\Phi$; and, on $\Delta$ with excluded middle,
the path type itself.

Throughout, $A$ is fibrant over $\Gamma$ and $a,b\in\Tm(\Gamma,A)$. For
$\rho\in\Gamma(X)$ an element $p\in\Path_A(a,b)(\rho)$ is a line
$p(f,j)\in A(W,\rho f)$ over the constant line at $\rho$, with
$p(e_0)=a\rho$ and $p(e_1)=b\rho$. The \emph{constant path} at $\rho$ is
$(f,j)\mapsto a(\rho f)$; it is in $\Path_A(a,a)(\rho)$.

\subsection{Swan's identity type}

\begin{hypothesis}[dependent conjunction]\label{hyp:dc}
Let $\psi\in\Phi(X)$ and let $\varphi_f\in\Phi(W)$, for $f:W\to X$
landing in $\psi$, be natural: $g^*\varphi_f=\varphi_{fg}$. Then the
sieve
\[
\psi\wedge\varphi=\{f\ \text{in}\ \psi\mid \varphi_f=\top\}
\]
is in $\Phi(X)$.
\end{hypothesis}

It is a sieve: if $f$ is in $\psi$ and $\varphi_f=\top$, then $fg$ is in
$\psi$ and $\varphi_{fg}=g^*\varphi_f=\top$. It is decidable, since
membership in $\psi$ and $\varphi_f=\top$ are. In both instances $\Phi(X)$
is the set of all decidable sieves, so the hypothesis holds. Its two
properties used below are: for $f$ in $\psi$,
$f^*(\psi\wedge\varphi)=\varphi_f$; and
$h^*(\psi\wedge\varphi)=h^*\psi\wedge(g\mapsto\varphi_{hg})$.

\begin{definition}[$\Id$]\label{def:id}
$\Id_A(a,b)(X,\rho)$ is the set of pairs $(p,\varphi)$ with
$p\in\Path_A(a,b)(\rho)$ and $\varphi\in\Phi(X)$ such that $p$ is
constant on $\varphi$: $p(f,j)=a(\rho f)$ for every $f$ landing in
$\varphi$ and every $j$. Restriction along $h$ is
$(p,\varphi)h=(ph,h^*\varphi)$. Reflexivity is
$\mathrm{refl}_a(\rho)=(\text{the constant path},\top)\in\Id_A(a,a)(\rho)$,
and $\pi:\Id_A(a,b)\to\Path_A(a,b)$ is the first projection.
\end{definition}

This is a type: $\varphi$ is coded by its characteristic map, and the
condition is a proposition. On $\varphi$ the two ends agree, $a(\rho f)=
b(\rho f)$.

\begin{proposition}[$\Id$ is fibrant]\label{prop:id-fib}
$\Id_A(a,b)$ has a fibration structure, natural in $\Gamma$ in the sense
of Theorem~\ref{thm:formers}.
\end{proposition}

\begin{proof}
Let $\gamma$ be a line of $\Gamma$ at $X$, $(p_0,\varphi_0)$ an element
over $\gamma(e)$ and $(\psi,u)$ a partial line ending at it, with
$u_f(j)=(p_{f,j},\varphi_{f,j})$. Let $p_1$ be the composite of the path
type (Theorem~\ref{thm:formers}) for the partial line $\pi u$ starting
at $p_0$. The family $f\mapsto\varphi_{f,e'}$ (the sieve of the end of
$u_f$) is natural on $\psi$; let $\varphi_1=\psi\wedge\varphi_{-,e'}$. The
composite is $(p_1,\varphi_1)$.

It is an element: for $f$ in $\varphi_1$, $p_1f=p_{f,e'}$ by the sides of
the path composition, and $p_{f,e'}$ is constant since
$\varphi_{f,e'}=\top$. Sides: for $f$ in $\psi$, $p_1f=p_{f,e'}$ and
$f^*\varphi_1=\varphi_{f,e'}$. Uniformity: that of $p_1$, and
$h^*\varphi_1$ is the dependent conjunction of the restricted data.
Naturality: the composite is formed pointwise from the composition of the
path type, which commutes with moves; the cartesian map of identity
types is that of path types with the identity on $\varphi$.
\end{proof}

\begin{theorem}[$J$, strictly]\label{thm:swan-j}
Let $\Gamma_a=\Gamma.A.\Id_A(a,q)$, with points $(\rho,b,w)$,
$w\in\Id_A(a\rho,b)$. Let $C$ be a type over $\Gamma_a$ with a composition
$\Comp_{e_0\to e_1}(C)$, and $d\in\Tm(\Gamma,C[a,\mathrm{refl}_a])$. There
is $J(d)\in\Tm(\Gamma_a,C)$ with
\[
J(d)(\rho,a\rho,\mathrm{refl}_a\rho)=d(\rho).
\]
\end{theorem}

\begin{proof}
Let $(\rho,b,(p,\varphi))\in\Gamma_a(X)$ and let $c=c_0$ be the connection
with $c(e_0,k)=e_0$ and $c(e_1,k)=k$. The line $\ell$ of $\Gamma_a$ at $X$ is
\[
\ell(f,j)=\bigl(\rho f,\ p(f,j),\ (p^{f,j},\ f^*\varphi\cup(j=e_0))\bigr),
\qquad p^{f,j}(g,k)=p(fg,c(jg,k)).
\]
$p^{f,j}$ is a path from $a$ to $p(f,j)$, constant on $f^*\varphi$ and on
$(j=e_0)$, so $\ell(f,j)$ is a point; $\ell$ is natural. By the connection
laws $\ell(e_0)=(\rho,a\rho,\mathrm{refl}_a\rho)$ and
$\ell(e_1)=(\rho,b,(p,\varphi))$. For $f$ in $\varphi$, $\ell f$ is the
constant line at $(\rho f,a\rho f,\mathrm{refl}_a\rho f)$. Define
$J(d)(\rho,b,(p,\varphi))$ as the composite along $\ell$ from $d(\rho)$,
with the side $\varphi$ given by the constant lines at $d(\rho f)$. The
side ends at $d(\rho)f$, as required.

$J(d)$ is natural because $\ell$, $\varphi$ and $d$ restrict and the
composition is uniform. At $(\rho,a\rho,\mathrm{refl}_a\rho)$, $\varphi=\top$
contains the identity, so by the sides of the composition the composite
is the end of the side, $d(\rho)$.
\end{proof}

The computation rule is an equality of elements, and it is stable under
substitution because both sides are natural. Only the composition from
$e_0$ to $e_1$ of the motive is used.

\begin{proposition}[$\Id$ and $\Path$]\label{prop:id-path}
$\pi:\Id_A(a,b)\to\Path_A(a,b)$ has an equivalence structure, natural in
$\Gamma$.
\end{proposition}

\begin{proof}
Let $q\in\Path_A(a,b)(\rho)$ and let $(\psi,v)$ be a partial element of
$\HFib_\pi$ at $(\rho,q)$, $v_f=((p_f,\varphi_f),\lambda_f)$ with
$\lambda_f$ a line of paths from $p_f$ to $qf$. Fill the partial line
$\lambda$ by $\mathrm{Fill}_{e_1}$ of the path type, from $q$ at $e_1$
(Proposition~\ref{prop:fill}); let $L$ be the filler and $p=L(e_0)$. Let
$\varphi=\psi\wedge(f\mapsto\varphi_f)$. Then $((p,\varphi),L)$ is an
element of the fibre: $p$ is constant on $\varphi$ because $pf=p_f$ for
$f$ in $\psi$. On $\psi$ it is $v_f$, since $f^*\varphi=\varphi_f$ and
$Lf=\lambda_f$; it is uniform because the filler and $\wedge$ are.
\end{proof}

So $\Id_A(a,b)$ and $\Path_A(a,b)$ are equivalent fibrant types, and $\Id$
has a strict $J$. The construction is generic: it uses
Hypothesis~\ref{hyp:dc}, unions, the points of $[1]$ and a connection, and
holds for posets and for $\Delta$.

\subsection{The path type on $\Delta$}

A composition $\Comp_{e\to e'}(A)$ is \emph{regular} if its composite
along a constant line of $\Gamma$, from $x$, with a partial line constant
in $j$, is $x$. With a regular composition of the motive, $J$ for the path
type, defined as in Theorem~\ref{thm:swan-j} without the cofibration
($\ell(f,j)=(\rho f,p(f,j),p^{f,j})$ and the empty side), computes
strictly: at $(\rho,a\rho,\text{constant path})$ the line $\ell$ is
constant.

\begin{lemma}\label{lem:reg-move}
A regular composition moves along a cartesian map to a regular
composition.
\end{lemma}

\begin{proof}
The maps $\sigma$ send constant lines and constant partial lines to
constant ones, and $\tau$ is injective.
\end{proof}

Regularity need not be preserved by the type formers. On $\Delta$ it does
not have to be: a type of the model is a code, whose classifier part is a
horn filling operation, and its composition is always that of $\El$,
moved along the code (Section~\ref{sec:delta}). So it suffices to make
the composition of $\El$ regular. Since a horn filling operation carries
no uniformity, it can be changed pointwise.

\begin{definition}[normal horn filling]
Horn data $h$ on $\Lam^k[n+1]$ at $\rho\in\Gamma([n+1])$ is
\emph{degenerate along $s_i$}, for $k\in\{i,i+1\}$, if $\rho=s_i^*\rho'$
and $h$ is the restriction of $s_i^*y$ for some $y\in A(\rho')$. A horn
filling operation is \emph{normal} if it extends such $h$ by $s_i^*y$.
\end{definition}

The element $y$ is the face of $h$ at the vertex of $\{i,i+1\}$ other
than $k$, so it is determined by $h$.

\begin{lemma}\label{lem:deg-unique}
If $h$ is degenerate along $s_i$ and along $s_l$, then
$s_i^*y=s_l^*y'$.
\end{lemma}

\begin{proof}
If $i\neq l$, then $k\in\{i,i+1\}\cap\{l,l+1\}$; say $l=i+1=k$. The faces
$i$ and $i+2$ are in the horn, so $y=d_is_{i+1}y'=s_id_iy'$ and
$y'=d_{i+2}s_iy=s_id_{i+1}y=s_id_iy'$. Then $s_i^*y=s_is_id_iy'$ and
$s_{i+1}^*y'=s_{i+1}s_id_iy'$, which are equal.
\end{proof}

With excluded middle every horn filling operation has a \emph{normalisation}:
on horn data degenerate along some $s_i$ (a proposition), the degenerate
extension; otherwise the given one. By Lemma~\ref{lem:deg-unique} the
choice of $i$ does not matter, and the normalisation is normal.

\begin{proposition}\label{prop:reg-delta}
If the horn filling operation is normal, the fillers of
Theorem~\ref{thm:hornbox}(2) are regular: for a filling problem over a
constant line at $\rho$, with sides constant in $j$ and end $x$, the
filler is the constant line at $x$. Hence both compositions are regular.
\end{proposition}

\begin{proof}
By induction on the dimension $N$ of the problem $b$. If $b$ is
degenerate along $s_i$, its filler is $s_i^*$ of the filler of $d_i^*b$,
a problem of the same kind, constant by induction. If a side contains the
identity, the filler is that side, constant. Otherwise the prism is
filled from its boundary and its lid at $e$: the faces are the fillers of
the face problems, constant by induction, and the lid is $x$. So the data
is the restriction of the constant prism $x\circ\mathrm{pr}$, and so is
the horn data of the first shuffle $P_k$. Since $\mathrm{pr}\circ P_k=s_k$,
this horn data is the restriction of $s_k^*x$ at $s_k^*\rho$; the missing
face of the horn is $d_k$ (the chain from $e_0$) or $d_{k+1}$ (the chain
from $e_1$), so the data is degenerate along $s_k$, and the normal
operation fills $P_k$ by $s_k^*x$, the constant prism along $P_k$. The
same holds for each next shuffle. So the filler is the constant line, and
its other end is $x$.
\end{proof}

\begin{theorem}\label{thm:path-j-delta}
On $\Delta$, take as the fibration of $\El$ Theorem~\ref{thm:hornbox}(2)
applied to the normalisation of $\El$'s horn filling operation. Then the
composition of every type $c^*\El$ is regular, and the path type has a
$J$ with strict computation rule for every motive of the model.
\end{theorem}

\begin{proof}
The composition of $c^*\El$ is that of $\El$ moved along $c$; it is
regular by Proposition~\ref{prop:reg-delta} and Lemma~\ref{lem:reg-move}.
Nothing else in Section~\ref{sec:delta} changes: Hypothesis~\ref{hyp:el}
only asks for a fibration of $\El$, and condition (C2) concerns the horn
filling operations, which are unchanged.
\end{proof}

Excluded middle is now used at two points, both in Section~\ref{sec:delta}:
deciding the degeneracy of filling problems, and of horn data. The
normalisation is not natural in $\Gamma$ (as for
Theorem~\ref{thm:hornbox}(2), the degeneracy of $\rho$ is not preserved
by substitutions), which is why it is applied once, to $\El$. On posets
the classifier is $\Fib$ itself: the composition of $\El$ at a code is the
code's own, so regularity would have to be preserved by the type formers;
there the strict identity type is Swan's.

%======================================================================
\appendix
\section{Correspondence with the formalisation}

\noindent Files of \texttt{seg/core}; planned ones, from \texttt{seg/model}, in italics.

\medskip
\begin{tabular}{p{3.6cm}p{0.9cm}p{9cm}}
section & layer & files\\\hline
\ref{sec:meta} Metatheory & L0 & V, VCode, Encode, Presheaf; FunExtUA (function extensionality from univalence) and, in All, the instance ModelUA of the whole development under univalence alone\\
\ref{sec:base} Base & L1 & CBase (with $\top$), CSieve\\
\ref{sec:rc} Relative contractibility & L2 & CPartial, CMove, CMoveLaws, CSub, COp, COpMove\\
\ref{sec:fib} Fibrations & L3 & CLine, CConn, CFillProb, CFill, CCart, CFibSub, CPath, CEquiv; Path and transport (L7): CPathProb, CPathComp, CTr, CTrOm, CTrB, CTrP, CTrE; $\Sigma$, $\Pi$: CSig, CPiE, CPi; naturality: CFillMove, CSigNat, CPiCart, CPiNat; Path: CPathFib, CPathCart, CPathNat\\
\ref{sec:univ} Small types, universe & L4 & CSmall, CUniv, CElC\\
\ref{sec:glue} Gluing, Theorem 7 & L5 & CExtL, CExt7Nat, CCof, CExt--CExt4, CExtEquiv--3, CExtComp1--4, CThm7, CExtFull, CLineSide, CMapLine, CHFibCart, CHFibLaws, CTrEq, CEnc, CGlueD, CGlueE, CGlueLaws, CGlueRC, CExtBig, CExtEq1--2, CExtFibU; \emph{GExtCmp, GExtSmall}\\
\ref{sec:ua} Univalence, $\Comp(1,\U)$ & L6, L7 & L6: CHSub, CUa, CGlueCart, CGlueCartNat, CUaD, CUaB, CUaN, CUaBN, CGlueUniq, CUaFull1, CExtFullCart, CUaSA, CUaSB, CUaOp; L7: CFibU1, CFibU2, CFibU3\\
\ref{sec:posets} Posets & L8 & Cat, CPosC, CPosBase, CPosInst, CPosEl1--5, CPosFinal; generic for $\Cl=\Fib$: CFibCl, CFibExt, CCartLaws, CSmallSub, CFillEnd, CExtFullFib, CTrFib, CYonCart, CSmallFull1--3, CFibFinal\\
\ref{sec:delta} Simplicial sets & L8 & CDeltaBase, CSimp, CDeltaNat, CDeltaFact, CDeltaCl, CBoxHorn, CDeltaRetr, CDGlue--CDNormE, CDeltaHornBox, CDeltaHornToFib, CDElHorn, CDeltaElFib, CDGlueCl, CDSmallCl, CDSmallFull(2), CDOpFull, CDOpYon, CDeltaFinal; generic for operations: COpCl, COpTr, COpLaws, CHornRetr\\
\ref{sec:id} Identity types & L3, L8 & Swan's $\Id$, generic: CId (type, refl, Hypothesis~\ref{hyp:dc} as DConj), CJ ($J$ from a contraction), CIdJ (Thm~\ref{thm:swan-j}), CIdFib (Prop~\ref{prop:id-fib}), CIdNat, CIdCart (its naturality), CIdEquiv (Prop~\ref{prop:id-path}), CIdEquivNat (its naturality); decidable sieves: CSieve.DConjS; instances: CPosId, CDeltaId. \emph{Section~9.2 not yet formalised}
\end{tabular}

\end{document}
